\chapter{Desarrollo}
\label{chap:desarrollo}

El proyecto principal consistió en implementar un \textbf{Blink LED} utilizando el LED 
integrado en la placa BluePill, el cual se encuentra conectado al pin \textbf{PC13}.

\section{Configuración de Periféricos}

Para controlar el LED, es necesario configurar el puerto C y el pin 13. El proceso 
implica los siguientes pasos:

\begin{enumerate}
    \item \textbf{Habilitación del reloj (Clock Gating)}: Se debe habilitar el reloj para 
    el puerto C en el registro \texttt{RCC\_APB2ENR}.
    \item \textbf{Configuración del pin}: Se configura el pin PC13 como salida en el 
    registro \texttt{GPIOC\_CRH}.
\end{enumerate}

\section{Archivos del Proyecto}

El firmware se compone de tres archivos fundamentales:

\begin{itemize}
    \item \texttt{main.c}: Contiene la lógica principal del programa y el bucle de parpadeo.
    \item \texttt{crt.s}: Archivo de \textit{startup} en ensamblador que inicializa el 
    puntero de pila y el vector de interrupciones.
    \item \texttt{linker.ld}: Script del enlazador que define la distribución de la 
    memoria Flash y RAM.
\end{itemize}

\section{Proceso de Compilación y Grabación}

Se utilizó un \texttt{Makefile} para automatizar la generación del binario. El flujo 
consiste en compilar los archivos fuente a objetos, enlazarlos para generar un archivo 
ELF, y finalmente extraer el binario (\texttt{.bin}).

Para la grabación se utilizó OpenOCD con el comando:
\begin{verbatim}
openocd -f interface/stlink.cfg -f target/stm32f1x.cfg
\end{verbatim}
seguido de los comandos de grabación en la consola de OpenOCD.
